% Imporditud: tools/import_eksperimendid.py
% Allikas: Eesti füüsikaolümpiaadi eksperimendi ülesannete
%   kogu (Rael Kalda, 2023), ülesanne Ü12, lahendus L12.
\setAuthor{EFO žürii}
\setRound{piirkonnavoor}
\setYear{2005}
\setNumber{P E1}
\setDifficulty{1}
\setTopic{dünaamika}
\setVahendid{mõõtejoonlaud, pall}
\setPunktid{10}
\setAllikas{Eksperimendi kogu Ü12, lahendus lk 38}
\setLahendusseis{toores}

\prob{Pall}
Kui ülestõstetud pall lahti lasta, muutub põrkel osa palli esialgsest potentsiaalsest energiast teisteks energia liikideks. Hinnata, kui suur osa algsest potentsiaalsest energiast muundub teisteks energia liikideks. Kas algsest potentsiaalsest energiast teisteks liikideks muutuva energia osa oleneb palli kukkumise kõrgusest?

\hint

\solu
Idee: kasutada potentsiaalse energia mõõduna raskusjõu poolt tehtavat tööd Ep = A = F s = mgh --- (2 p); mõõtmise planeerimine (teha mõõtmisi kahel oluliselt erineval kõrgusel, kordusmõõtmised, keskmise leidmine) --- (3 p); muundunud energiaosa hindamine põrke kõrguse järgi $\Delta$E = (halg $-$ hlopp)/halg --- (3 p); tulemuste võrdlemine --- (1 p); järeldus (ei olene) --- (1 p).
\probend
